\begin{statement}{E4}{Proved}{Degree-one cohomological flatness over a curve}
For a proper surjective holomorphic map $g:Y\to C$ from a connected complex
manifold to a nonsingular curve with connected fibres,
$R^1g_*\mathcal O_Y$ is locally free and
$H^0(Y_c,\mathcal O_{Y_c})=\mathbb C$ for every scheme-theoretic fibre.
\end{statement}
\begin{proof}
We rewrite \cite[Theorem labelled \texttt{raynaud}]{Viaclovsky2026Combs}.
Work over a disc with coordinate
$t$. Multiplication by $t$ is injective on $\mathcal O_Y$, so the map
is flat. For a proper map $h:Z\to\Delta_u$ with connected reduced
central fibre, $h_*\mathcal O_Z=\mathcal O_{\Delta_u}$, and $u$ a
nonzero divisor, the exact sequence for multiplication by $u$ shows
that $R^1h_*\mathcal O_Z$ has no $u$-torsion: the preceding map
$\mathcal O_{\Delta_u,0}\to H^0(Z_0,\mathcal O_{Z_0})=\mathbb C$
is surjective.

Resolve the divisor $g^{-1}(0)$ by a proper modification
$\mu:\widetilde Y\to Y$ with nonsingular source, so its pullback has
simple normal-crossing support and local equation
$t=z_1^{m_1}\cdots z_r^{m_r}$. Choose $M$ divisible by every
multiplicity, make $t=u^M$, and normalize. Denote the resulting
normal space by $Z$ and its map to $\Delta_u$ by $h$. Its central
fibre is reduced. Indeed, putting $e_i=M/m_i$, each local normal
component is the quotient of a polydisc by
\[
 G_0=\{(\zeta_i):\zeta_i^{e_i}=1,\ \prod\zeta_i=1\},
 \quad z_i=x_i^{e_i},\quad u=\zeta\prod x_i,
\]
where the appropriate constant root of unity $\zeta$ specifies the
chosen component. This includes the distinct components when the
exponents have a common divisor. This finite normal model is
birational onto that component.
For its power-series ring $R$, averaging gives
$R^{G_0}/uR^{G_0}=(R/uR)^{G_0}$, which is reduced. The Stein factor
of $h$ has generic degree one, since the general fibre is connected;
normality makes it $\Delta_u$. Thus $Z_0$ is connected and the
preceding argument gives local freeness of $R^1h_*\mathcal O_Z$.

The deck group $G=\mu_M$ has quotient $Z/G=\widetilde Y$.
To verify this, the invariant subring of the integral normalization
of a local base-change algebra lies in the invariant fraction field
and is integral over the normal ring of $\widetilde Y$; hence it is
that ring. Exactness of finite pushforward and of averaging gives
\[
 R^1(g\mu)_*\mathcal O_{\widetilde Y}
   =\bigl(b_*R^1h_*\mathcal O_Z\bigr)^G,
 \qquad b(u)=u^M.
\]
The right side is a direct summand of a locally free sheaf, since
$b$ is finite flat. The degree-one Leray sequence for $g\mu$ injects
$R^1g_*\mathcal O_Y$ into this sheaf, because
$\mu_*\mathcal O_{\widetilde Y}=\mathcal O_Y$. It is torsion-free,
and therefore locally free on the nonsingular curve. Finally pushing
forward $0\to\mathcal O_Y\xrightarrow{t}\mathcal O_Y\to
\mathcal O_{Y_0}\to0$ and using the absence of $t$-torsion proves
$H^0(Y_0,\mathcal O_{Y_0})=\mathbb C$.
\end{proof}


% E5 from part_E.tex

